Let us first show that the morphism of (ii) is an immersion. By definition, $\operatorname{int}(n_w)^{-1}$ induces a closed immersion of $B'{}^u_w$ into $B'{}^u$, hence the morphism
\[
(u,b)\mto n_w^{-1}un_wb
\]
induces a closed immersion
\[
B'{}^u_w\underset S\times B\to\Omega_{R_+}.
\]
This implies immediately that the morphism of (ii) induces a closed immersion of the first side into the open subset $n_w\Omega_{R_+}$. To prove (i), it suffices to see that
\[
B'(S)N_w(S)B(S)=B'{}^u_w(S)N_w(S)B^u(S).
\]
Now, if $\alpha\in R$, one has $\operatorname{int}(n_w)U_\alpha(S)=U_{w(\alpha)}(S)$, hence if $w^{-1}(\alpha)\in R_+$,
\[
\begin{aligned}
U_\alpha(S)N_w(S)B(S)&=U_\alpha(S)n_wT(S)B^u(S)\\
&=n_wU_{w^{-1}(\alpha)}(S)T(S)B^u(S)\\
&=n_wB(S)=N_w(S)B^u(S).
\end{aligned}
\]
This implies at once, in view of the definition of $R_-^w$, the desired assertion.
\begin{theorem}{5.7.4}\label{XXII.5.7.4}
Let $S$ be a scheme, $(G,T,M,R)$ a split $S$-group, $B$ the Borel subgroup defined by the system of positive roots $R_+$, $B'$ the Borel subgroup defined by $R_-=-R_+$.
\begin{enumeratei}
\item (Bruhat's theorem) The schemes $B'{}^u_w\cdot N_w\cdot B$ form, as $w$ ranges over $W$, a partition of the underlying set of $G$.
\item For each $w\in W$, let $n_w$ be a representative of $w$ in $\uNorm_G(T)(S)$ (3.8); then the open subsets $n_w\Omega=n_wB'{}^u\cdot T\cdot B^u$ form, as $w$ ranges over $W$, a covering of $G$.
\end{enumeratei}
\end{theorem}
The two assertions are indeed verified on the geometric fibers, where one concludes by 5.7.1 and 5.7.3.
\begin{remark}{5.7.5}\label{XXII.5.7.5}
(i) implies that if $S$ is the spectrum of a field, $G(S)$ is the disjoint union of the $B'{}^u_w(S)\cdot T(S)\cdot B^u(S)$. The corresponding assertion for arbitrary $S$ (even local or artinian) is clearly false. Note however that (ii) implies that if $S$ is local, $G(S)$ is indeed the union of the $n_w\Omega(S)$. In fact:
% typo?: kept as printed: no representative n_w appears in the disjoint union.
\end{remark}
\begin{corollary}{5.7.6}\label{XXII.5.7.6}
Let $\Delta$ be a system of simple roots of the split group $G$ over the local scheme $S$.
\begin{enumeratei}
\item Then $G(S)$ is generated by $T(S)$ and the $U_\alpha(S)$, $\alpha\in\Delta\cup-\Delta$.
\item If $G$ is simply connected (4.3.3), $G(S)$ is already generated by the $U_\alpha(S)$, $\alpha\in\Delta\cup-\Delta$.
\end{enumeratei}
\end{corollary}
Indeed, let $H$ be the subgroup of $G(S)$ generated by the $U_\alpha(S)$, $\alpha\in\Delta\cup-\Delta$. Note first that $H$ contains a representative of each $s_\alpha$ ($\alpha\in\Delta$) in $\uNorm_G(T)(S)$ (Exp.~XX 3.1), hence a representative $n_w$ of each $w\in W$.

Since every $\alpha\in R$ is written $w(\alpha_0)$ with $w\in W$, $\alpha_0\in\Delta$, one has
\[
U_\alpha(S)=\operatorname{int}(n_w)U_{\alpha_0}(S)\subset H.
\]
The subgroup generated by $H$ and $T(S)$ therefore contains $\Omega(S)$ and is therefore all of $G(S)$, by the remark made earlier.

If now $G$ is simply connected, let us prove that $H\supset T(S)$. By Exp.~XX 2.7, $H$ contains $\alpha^*(\mathbf G_m(S))$ for every $\alpha\in\Delta$, and it suffices to apply 4.3.8.
\begin{remark}{5.7.6.1}\label{XXII.5.7.6.1}
Instead of taking, for each $\alpha\in\Delta$, $U_\alpha(S)$ and $U_{-\alpha}(S)$, one may content oneself with taking $U_\alpha(S)$ and a representative $w_\alpha$ of the reflection $s_\alpha$, or $U_\alpha(S)$ and a section of $U_{-\alpha}^\times$, $\ldots$.
\end{remark}
\begin{corollary}{5.7.7}\label{XXII.5.7.7}
If $G$ has semisimple rank $1$, choose a $u_\alpha\in U_\alpha^\times(S)$. Then $\Omega$ and $u_\alpha\Omega$ form a covering of $G$.
\end{corollary}
Indeed, if $u_{-\alpha}$ is the section of $U_{-\alpha}$ paired with $u_\alpha$ (\cf 1.3), one has, by 5.7.4 (ii),
\[
G=\Omega\cup u_{-\alpha}^{-1}u_\alpha u_{-\alpha}^{-1}\Omega,
\]
whence
\[
G=u_{-\alpha}G=u_{-\alpha}\Omega\cup u_\alpha u_{-\alpha}^{-1}\Omega=\Omega\cup u_\alpha\Omega.
\]
\begin{corollary}{5.7.8}\label{XXII.5.7.8}
Let $S$ be a scheme, $G$ a reductive $S$-group. Then $G$ is essentially free over $S$ (Exp.~VIII 6.1).{\renewcommand{\thefootnote}{(47)}\nde{See also the additions made in VIB, 6.2.1 to 6.2.6 and 6.5.2 to 6.5.5.}}
\end{corollary}
Indeed, the assertion is local for the (fpqc) topology; one may assume $G$ split. Then $G$ admits a covering by open subsets isomorphic to $\mathbf G_{a,S}^N\times_S\mathbf G_{m,S}^n$, hence essentially free.
\begin{lemma}{5.7.9}\label{XXII.5.7.9}
Under the conditions of 5.7.4, let $\alpha$ be a simple root of $R_+$ and $u_\alpha\in U_\alpha^\times(S)$. For every $v\in U_{-\alpha}(S)$, one has
\[
\Omega\cdot v\subset\Omega\cup u_\alpha\cdot\Omega.
\]
\end{lemma}
One must compare two open subsets of $G$; it suffices to do so when $S$ is the spectrum of a field $k$. One must therefore prove
\[
\Omega(k)v\subset\Omega(k)\cup u_\alpha\Omega(k).
\]
Now
\[
\begin{aligned}
\Omega(k)v&=B'{}^u(k)T(k)B^u(k)v\\
&=U_{-\widehat\alpha}(k)U_{-\alpha}(k)T(k)U_\alpha(k)U_{\widehat\alpha}(k)v\\
&\subset U_{-\widehat\alpha}(k)Z_\alpha(k)U_{\widehat\alpha}(k)v.
\end{aligned}
\]
(One uses the decomposition of 5.6.8). Applying now 5.6.8 (iii) and using 5.7.7 for the group $Z_\alpha$, one obtains
\[
\begin{aligned}
\Omega(k)v&\subset U_{-\widehat\alpha}(k)Z_\alpha(k)vU_{\widehat\alpha}(k)\subset U_{-\widehat\alpha}(k)Z_\alpha(k)U_{\widehat\alpha}(k)\\
&\subset U_{-\widehat\alpha}(k)U_{-\alpha}(k)T(k)U_\alpha(k)U_{\widehat\alpha}(k)\\
&\qquad\cup U_{-\widehat\alpha}(k)u_\alpha U_{-\alpha}(k)T(k)U_\alpha(k)U_{\widehat\alpha}(k).
\end{aligned}
\]
Using 5.6.8 (iii) again (for $R_-$ instead of $R_+$), one obtains the result.
\begin{proposition}{5.7.10}\label{XXII.5.7.10}
Under the conditions of 5.7.4, choose for each simple root $\alpha$ a $u_\alpha\in U_\alpha^\times(S)$. Let $U_1$ be the submonoid of $B^u(S)$ generated by the $u_\alpha$. The open subsets $u\Omega$, for $u\in U_1$, form a covering of $G$.
\end{proposition}
Once again, one may assume that $S$ is the spectrum of a field $k$; by 5.7.6, it suffices to prove that $\bigcup_{u\in U_1}u\Omega(k)$ is stable under right multiplication by $T(k)$, $U_\alpha(k)$, $U_{-\alpha}(k)$ (for $\alpha$ simple). In the first two cases, this is trivial. In the last, this follows from the lemma.
\begin{remark}{5.7.11}\label{XXII.5.7.11}
Let us point out a particular case of 5.7.2. If $w=s_\alpha$ is the reflection with respect to the simple root $\alpha$, then
\[
R_-\cap s_\alpha(R_-)=R_--\{-\alpha\}
\]
(Exp.~XXI 3.3.1), and, in the notations of 5.6.8, one therefore has
\[
B'{}^u_{s_\alpha}=U_{-\widehat\alpha}.
\]
\end{remark}
\begin{remark}{5.7.12}\label{XXII.5.7.12}
In fact, the proof of 5.7.10 gives at once the following statement: under the conditions of 5.7.10, let $\Gamma$ be a submonoid of $G(S)$; for the open subsets $g\Omega$ ($g\in\Gamma$) to form a covering of $G$, it is necessary and sufficient that for every $s\in S$ and every simple root $\alpha$, one have
\[
(u_\alpha)_{\overline s}B'{}^u(\overline s)\subset\Gamma\cdot B'{}^u(\overline s)\cdot T(\overline s)\cdot B^u(\overline s).
\]
\end{remark}
\begin{remark}{5.7.13}\label{XXII.5.7.13}
By 5.5.5 (iii), arguing as in 5.7.1, one obtains at once the following variant of 5.7.4: let $(G,T,M,R)$ be a split $S$-group, $B$ and $B'$ two Borel subgroups of $G$ containing $T$; for every $w\in W$, the sheaf $B'\cdot N_w\cdot B$ is representable by a subscheme of $G$; these subschemes form, for $w\in W$, a partition of the underlying set of $G$. One may also give the analogue of 5.7.3 (ii): one must put
\[
B'{}^u_w=B'{}^u\cap\operatorname{int}(n_w)\widetilde B^u,
\]
where $\widetilde B$ is the Borel subgroup ``opposite'' to $B$ relative to $T$ (\cf 5.9.2).
\end{remark}
\begin{proposition}{5.7.14}\label{XXII.5.7.14}
Let $S$ be a scheme, $G$ a reductive $S$-group, and
\[
\Ad:G\to\GL_{\OS}(\mathfrak g)
\]
its adjoint representation. Then $\Ker(\Ad)=\uCentr(G)$ (in other words, the canonical homomorphism deduced from $\Ad$ on passing to the quotient:
\[
\overline{\Ad}:G/\uCentr(G)=\operatorname{ad}(G)\to\GL_{\OS}(\mathfrak g)
\]
is a monomorphism).{\renewcommand{\thefootnote}{(48)}\nde{It is even a closed immersion, by Exp.~XVI 1.5 (a).}}
\end{proposition}
One may assume $G$ split. Choose on $\Gamma_0(R)$ a total order structure compatible with the group structure and let $R_+$ be the set of positive roots. By 5.7.4 (ii) and 4.1.6, it suffices to prove that if $n_w$ is a representative of the element $w$ of $W$, if $u\in U(S)$, $t\in T(S)$, $v\in U_-(S)$, and if $\Ad(n_wvtu)=\mathrm{id}$, then $w=e$, $v=e$, $u=e$. For each $m\in R\cup\{0\}$, put
\[
\mathfrak g^{>m}=\coprod_{n>m}\mathfrak g^n,\qquad\mathfrak g^{<m}=\coprod_{n<m}\mathfrak g^n.
\]
Let $X\in\Gamma(S,\mathfrak g^m)$; write $\Ad(tu)X=\Ad(v^{-1}n_w^{-1})X$. Now
\[
\begin{aligned}
\Ad(t)\Ad(u)X-m(t)X&\in\Gamma(S,\mathfrak g^{>m}),\\
\Ad(v^{-1}n_w^{-1})X-\Ad(n_w^{-1})X&\in\Gamma(S,\mathfrak g^{<w^{-1}(m)}).
\end{aligned}
\]
If $w\ne e$, there exists an $\alpha\in R$ such that $w^{-1}(\alpha)<\alpha$, and putting $m=\alpha$, one obtains a contradiction, for
\[
\begin{gathered}
\Ad(tu)X\in\Gamma(S,\mathfrak g^\alpha+\mathfrak g^{>\alpha})\cap
\Gamma(S,\mathfrak g^{w^{-1}(\alpha)}+\mathfrak g^{<w^{-1}(\alpha)})=0.
\end{gathered}
\]
One therefore has $w=e$, and one may choose $n_w=e$; one then has
\[
\Ad(v^{-1})X-X\in\Gamma(S,\mathfrak g^{<m}\cap(\mathfrak g^m+\mathfrak g^{>m}))=0,
\]
whence $\Ad(v)X=X$ for every $X\in\Gamma(S,\mathfrak g^m)$, hence $\Ad(v)=\mathrm{id}$. Likewise $\Ad(u)=\mathrm{id}$. One then concludes by 5.6.2 bis.

\sgasection{5.8}{Schemes associated with a reductive group}
\begin{theorem}{5.8.1}\label{XXII.5.8.1}
Let $S$ be a scheme, $G$ a reductive $S$-group. Let $\mathscr H$ be the functor of subgroups of type (R) of $G$: for every $S'\to S$, $\mathscr H(S')$ is the set of subgroups of type (R) of $G_{S'}$ (\cf 5.2.1). Then $\mathscr H$ is representable by a quasi-projective $S$-scheme, of finite presentation over $S$.
\end{theorem}
{\renewcommand{\thefootnote}{(49)}\nde{The original has been expanded in what follows.}}%
Let $G'=G/\uCentr(G)$ be the adjoint group of $G$ (4.3.6) and $u$ the morphism $G\to G'$. By Exp.~XII 7.12, the map $H'\mto u^{-1}(H')$ establishes a bijection between the set of subgroups of type (R) of $G'$ and of $G$ (and this remains valid after every base change). Hence, replacing $G$ by $G'$, one may assume that $G$ is adjoint. Consider then the morphism
\[
u:\mathscr H\to\underline{\operatorname{Grass}}(\mathfrak g)
\]
which associates with each subgroup of type (R) its Lie algebra (which is a submodule locally a direct summand of $\mathfrak g${\renewcommand{\thefootnote}{(50)}\nde{\cf Exp.~II 4.11.8.}}). Then $u$ is a monomorphism by 5.3.3. It suffices to prove that it is representable by an immersion of finite presentation, in other words to prove the following assertion: for every $S_1\to S$, given a submodule $\mathfrak h$ locally a direct summand of $\mathfrak g_{S_1}$, the $S'\to S_1$ such that $\mathfrak h_{S'}$ is the Lie algebra of a subgroup of type (R) of $G_{S'}$ are exactly those which factor through a certain subscheme $\Sigma$ of finite presentation of $S_1$. Replacing $S_1$ by $S$, one reduces to $S_1=S$, and one may moreover assume $S$ affine; then there exists an affine noetherian scheme $S_0$ such that $G$ (resp. $\mathfrak h$) comes by base change from an adjoint reductive $S_0$-group $G_0$ (resp. a submodule $\mathfrak h_0$ locally a direct summand of $\mathfrak g_0=\uLie(G_0/S_0)$). It suffices to show that there exists a subscheme $\Sigma_0$ of $S_0$ having the required properties (for then one will have $\Sigma=\Sigma_0\times_{S_0}S$). Replacing $S$ by $S_0$, one may therefore assume $S$ affine and noetherian (note that then every subscheme of $S$ is of finite presentation over $S$). Finally, replacing $S$ by a sufficiently small open subset, one may assume that $\mathfrak g$ is free of rank $n$ and that $\mathfrak h$ is a direct summand, free of rank $r$.

One must first express that $\mathfrak h_{S'}$ is a Lie subalgebra of $\mathfrak g_{S'}$, i.e. that the morphism induced by the Lie bracket:
\[
\varphi:\mathfrak h\otimes\mathfrak h\xrightarrow{[\ ,\ ]}\mathfrak g
\]
factors through $\mathfrak h$. If $(e_1,\ldots,e_n)$ is a basis of $\mathfrak g$ such that $(e_1,\ldots,e_r)$ is a basis of $\mathfrak h$, then $\varphi$ is given by sections $a_k^{ij}$ of $\OS$ (where $i,j=1,\ldots,r$ and $k=1,\ldots,n$), and the preceding condition is equivalent to saying that $S'\to S$ factors through the closed subscheme of $S$ defined by the equations $a_k^{ij}=0$ for $k=r+1,\ldots,n$ and $i,j=1,\ldots,r$. Replace $S$ by this closed subscheme.

Then, by 5.3.0, $N=\uNorm_G(\mathfrak h)$ is a closed group subscheme of $G$, of finite presentation over $G$. One must now express (\cf 5.3.1) that $N_{S'}$ is smooth at every point of the identity section, of relative dimension $r=\operatorname{rank}(\mathfrak h)$, and that the inclusion of $\mathfrak h_{S'}$ into $\uLie(N_{S'}/S')$ is an equality.

{\renewcommand{\thefootnote}{(51)}\nde{The original next indicated that, denoting by $\mathfrak n_{N/G}$ the inverse image of the conormal sheaf $\mathcal N_{N/G}$ by the identity section $S\to N$, the condition that $\mathfrak n_{N_{S'}/G_{S'}}\to\omega^1_{G_{S'}/S'}$ be universally injective is equivalent to the fact that $S'\to S$ factors through a certain open subscheme of $S$. We have not succeeded in justifying this point, owing to the fact that the formation of $\mathcal N_{N/G}$ does not commute with base change, and we have replaced this argument by the one which follows, indicated by O.~Gabber.}}%
Since $N$ is affine over $S$ (being closed in $G$), the identity section $\varepsilon:S\to N$ is a closed immersion, so $\varepsilon(S)$ is defined by a quasi-coherent ideal $\mathcal J$ of $\mathcal A(N)$. Note that $\mathfrak n_{S/N}=\mathcal J/\mathcal J^2$ is identified with $\varepsilon^*(\Omega^1_{N/S})$, hence its formation commutes with every base change $S'\to S$. By equivalence (c$'$) $\Leftrightarrow$ (a) in EGA IV$_4$, 17.12.1 (applied to $f:N\to S$ and $j=\varepsilon$), $N_{S'}\to S'$ is smooth, of relative dimension $r$, at every point of $\varepsilon(S')$ if and only if $\mathfrak n_{S'/N'}=\mathfrak n_{S/N}\otimes_{\OS}\OO_{S'}$ is locally free of rank $r$ and the morphism $\varphi_n(S'):\operatorname{Sym}^n(\mathfrak n_{S'/N'})\to\mathcal J_{S'}^n/\mathcal J_{S'}^{n+1}$ is an isomorphism for every $n\ge1$. Put $K_n(S')=\Ker\varphi_n(S')$. By TDTE I, Lemma 3.6, $\mathfrak n_{S/N}\otimes_{\OS}\OO_{S'}$ is locally free of rank $r$ if and only if $S'\to S$ factors through a certain subscheme $Z$ of $S$. Replacing $S$ by $Z$, one may therefore assume that $\mathfrak n_{S/N}=\mathcal J/\mathcal J^2$ is locally free of rank $r$. Then, for every $S'\to S$, one has $(\mathcal J^2/\mathcal J^3)\otimes_{\OS}\OO_{S'}=\mathcal J_{S'}^2/\mathcal J_{S'}^3$ and hence $K_2(S')=K_2(S)\otimes_{\OS}\OO_{S'}$. It follows that $\varphi_2(S')$ is an isomorphism if and only if $S'\to S$ factors through the closed subscheme $S_2$ of $S$ defined by the ideal generated by the image of $\operatorname{Sym}^2(\mathfrak n_{S/N})^*\otimes K_2(S)$ in $\OS$. Then, over $S_2$, $\mathcal J^2/\mathcal J^3$ is isomorphic to $\operatorname{Sym}^2(\mathfrak n_{S/N})$, hence locally free, and the same argument shows that $\varphi_3(S')$ is an isomorphism if and only if $S'\to S$ factors through a certain closed subscheme $S_3$ of $S$, etc. One thus obtains that $N_{S'}\to S'$ is smooth, of relative dimension $r$, at every point of $\varepsilon(S')$ if and only if $S'\to S$ factors through the closed subscheme $Z$ which is the intersection of the $S_n$. But then, for every $S'\to Z$, $\uLie(N_{S'}/S')$ is locally a direct summand of rank $r$ of $\mathfrak g_{S'}$, and hence the inclusion $\mathfrak h_{S'}\subset\uLie(N_{S'}/S')$ is an equality. One then puts $H=N^0$.

Replacing $S$ by $Z$, it only remains to express that $H_{s'}$ has the same reductive rank as $G_{s'}$ at every point $s'\in S'$, or, equivalently, that $H_s$ has the same reductive rank as $G_s$ at every point $s$ of the (set-theoretic) image of $S'$ in $S$. Now this condition defines an open subset of $S$ (Exp.~XIX 6.2).
\begin{remark*}
In general, the scheme $\mathscr H$ is not smooth over $S$. It is however if $6$ is invertible on $S$, or if there exists a prime number $p$ such that $p\cdot1_S=0$ (i.e. if $S$ has characteristic $p>0$).
\end{remark*}
\begin{corollary}{5.8.2}\label{XXII.5.8.2}
Let $S$ be a scheme, $G$ a reductive $S$-group, $H$ a subgroup of type (R) of $G$. (Recall (5.3.10) that $\uNorm_G(H)$ is representable by a closed group subscheme of $G$, smooth over $S$).

Then the quotient sheaf $G/\uNorm_G(H)$ is representable by a quasi-projective $S$-scheme, smooth and of finite presentation over $S$ (which is in fact an open subset of $\mathscr H$).
\end{corollary}
Indeed, consider the morphism
\[
f:G\to\mathscr H,
\]
defined set-theoretically by $f(g)=\operatorname{int}(g)H$. By 5.3.9, this morphism is smooth and of finite presentation, hence open. Let $V=f(G)$ be endowed with its structure of open subscheme of $\mathscr H$. The morphism $f:G\to V$ is covering and has kernel $\uNorm_G(H)$, which proves that $G/\uNorm_G(H)$ is representable by $V$ (\cf Exp.~IV 4.6.5).
\begin{corollary}{5.8.3}\label{XXII.5.8.3}
Let $S$ be a scheme, $G$ a reductive $S$-group. Consider the functors $\underline{\Tor}(G)$, $\underline{\operatorname{Bor}}(G)$, $\underline{\operatorname{Kil}}(G)$ defined by
\[
\begin{aligned}
\underline{\Tor}(G)(S')&=\{\text{maximal tori of }G_{S'}\},\\
\underline{\operatorname{Bor}}(G)(S')&=\{\text{Borel subgroups of }G_{S'}\},\\
\underline{\operatorname{Kil}}(G)(S')&=\{\text{Killing pairs of }G_{S'}\text{ (\cf 5.3.13)}\}.
\end{aligned}
\]
\begin{enumeratei}
\item $\underline{\Tor}(G)$, $\underline{\operatorname{Bor}}(G)$, $\underline{\operatorname{Kil}}(G)$ are representable by smooth $S$-schemes of finite presentation, with integral geometric fibers, and respectively affine, projective, affine over $S$.
\item The canonical morphism $\underline{\operatorname{Kil}}(G)\to\underline{\Tor}(G)$ (resp. $\underline{\operatorname{Kil}}(G)\to\underline{\operatorname{Bor}}(G)$) is finite \'etale surjective (resp. affine smooth surjective).
\item Let $T$ be a maximal torus of $G$ (resp. $B$ a Borel subgroup of $G$, resp. $B\supset T$ a Killing pair of $G$). The morphism
\[
G\to\underline{\Tor}(G),\quad\text{resp. }G\to\underline{\operatorname{Bor}}(G),\quad\text{resp. }G\to\underline{\operatorname{Kil}}(G)
\]
defined by
\[
g\mto\operatorname{int}(g)T,\quad\text{resp. }g\mto\operatorname{int}(g)B,\quad\text{resp. }g\mto(\operatorname{int}(g)B,\operatorname{int}(g)T)
\]
induces an isomorphism
\[
\begin{gathered}
G/\uNorm_G(T)\isomto\underline{\Tor}(G),\qquad\text{resp. }G/B\isomto\underline{\operatorname{Bor}}(G),\\
\text{resp. }G/T\isomto\underline{\operatorname{Kil}}(G).
\end{gathered}
\]
\end{enumeratei}
\end{corollary}
